\chapter{Sobolev and isoperimetric inequalities}\label{sec:submanifold-inequalities}
The area functional connects intrinsic analysis with the extrinsic mean curvature. Here $\vec H$ is the trace of the second fundamental form, rather than its average.

\begin{theorem}[Michael--Simon inequality in the smooth case]\label{thm:michael-simon}
Let $\Sigma^n\subset\mathbb R^N$ be a smooth embedded submanifold, with $n\geq2$, and let $f\in C_c^1(\Sigma\setminus\partial\Sigma)$ be nonnegative. There is a dimensional constant $C(n,N)$ such that
\[
 \left(\int_\Sigma f^{n/(n-1)}\,dA\right)^{(n-1)/n}
 \leq C(n,N)\int_\Sigma\bigl(|\nabla_\Sigma f|+|\vec H|f\bigr)\,dA.
\]
The compact support is intrinsic and avoids the boundary. This is the smooth-submanifold specialization of \cite{MS73}, Example 1.6 and Theorem 2.1. We do not claim a sharp value of $C(n,N)$. The proof remains external.
\end{theorem}

\begin{theorem}[Sharp isoperimetry in codimension at most two]\label{thm:brendle-isoperimetry}
Let $\Sigma^n$ be a smooth compact embedded $n$-dimensional submanifold of $\mathbb R^{n+2}$ with smooth boundary, where $n\geq2$. If $\omega_n$ is the volume of the unit ball in $\mathbb R^n$, then
\[
 \mathcal H^{n-1}(\partial\Sigma)+\int_\Sigma |\vec H|\,dA
 \geq n\omega_n^{1/n}\mathcal H^n(\Sigma)^{(n-1)/n}.
\]
In particular, the curvature term vanishes when $\Sigma$ is minimal. The same inequality applies to submanifolds of $\mathbb R^{n+1}$ by the isometric inclusion into $\mathbb R^{n+2}$. Source: \cite{Brendle19}, Corollary 2 in the supplied August 29, 2019 version. The proof remains external. This sharp constant is not asserted here in higher codimension.
\end{theorem}

\begin{example}[The flat ball fixes the normalization]\label{ex:flat-ball-isoperimetry}
For the radius-$R$ ball in an affine $n$-plane, $\vec H=0$, its $n$-volume is $\omega_nR^n$, and its boundary area is $n\omega_nR^{n-1}$. Substitution gives equality in Theorem \ref{thm:brendle-isoperimetry}. Both sides scale as $R^{n-1}$, as does the curvature integral because mean curvature scales as inverse length.
\end{example}
